% !PRINTED 458
\weberpage{458}
Für die Punkte $P_6$ ist dies unmittelbar einzusehen, wenn man
[nach \S.\,119, (1)] $x_1 = x_2 = x_3 = 1$ setzt, wodurch
$f(x_1, x_2, x_3) = 3$ wird, also nicht verschwindet. Um dasselbe für
die $P_8$ nachzuweisen, setzt man nach \S.\,118 (10):
\[
x_1 = \alpha\gamma\varepsilon^4, \quad x_2 = \beta\gamma\varepsilon,
\quad x_3 = \alpha\beta\varepsilon^2,
\]
wodurch, da $\alpha\beta\gamma = \frac{1}{7}$ ist [\S.\,115, (14)],
\[
f(x_1, x_2, x_3) = \frac{1}{7}(\alpha^3\gamma^2 + \beta^3\alpha^2 + \gamma^3\beta^2)
\]
wird. Setzt man darin [nach \S.\,115, (8)]:
\[
\alpha^3 = \alpha\beta\gamma - \alpha^2, \quad \beta^3 = \alpha\beta\gamma - \beta^2,
\quad \gamma^3 = \alpha\beta\gamma - \gamma^2,
\]
so erhält man
\weq{6}{f(x_1, x_2, x_3) = \tfrac{1}{7}\alpha\beta\gamma(\alpha^2+\beta^2+\gamma^2) - \tfrac{1}{7}(\alpha^2\gamma^2+\beta^2\alpha^2+\gamma^2\beta^2).}
Aus \S.\,115, (6), (7) aber folgt
\[
\alpha^2 + \beta^2 + \gamma^2 = 1,
\]
\[
\alpha^2\gamma^2 + \beta^2\alpha^2 + \gamma^2\beta^2 = -\alpha\beta\gamma(\alpha+\beta+\gamma) = \frac{1}{7},
\]
und daher ist
\[
f(x_1, x_2, x_3) = \frac{1}{42}
\]
von Null verschieden.

Die sämmtlichen Schnittpunkte der 21 Hauptaxen unter einander sind
die Pole $P_6$ und $P_8$, und folglich schneiden sich niemals zwei
Hauptaxen auf der Grundcurve. Die Schnittpunkte der Hauptaxen mit der
Grundcurve können also nur vierzählige oder zweizählige Pole sein. Um
darüber zu entscheiden, stellen wir folgende Erwägung an. Wenn unter
den Schnittpunkten der Hauptaxe $A$ mit der Grundcurve ein $P_2$
vorkommt, so sind alle vier Schnittpunkte von einander verschieden und
sind alle zweizählige Pole. Wenn aber ein $P_4$ auf $A$ und $f$ liegt,
so liegt auch der zweite auf $A$ liegende $P_4$ auf $f$, und es kann
keinen anderen Schnittpunkt von $A$ und $f$ geben, weil ein solcher
ein $P_2$ wäre und vier weitere $P_2$ zur Folge hätte. Die vier
Schnittpunkte von $A$ und $f$ müssen also paarweise zusammenfallen und
$A$ ist eine Doppeltangente. Nun haben wir aber gesehen, dass die 28
Linien $T$ Doppeltangenten sind, und eine Curve vierter Ordnung kann
nicht mehr als 28 Doppeltangenten haben. Daher ist die Annahme
unzulässig, und es folgt, dass die Schnittpunkte der Grundcurve mit
den Hauptaxen zweizählige Pole sind.

Bezeichnen wir also die besonderen zweizähligen Pole, die auf der
Curve $f$ liegen, mit $\bar P_2$, so haben wir folgende ausgezeichnete
Punktsysteme:
%CONT>>
